\subsubsection{Subproceso 3.B: Confirmación y Registro Oficial}

\begin{itemize}
    \item \textbf{Regla de aprobación automática (TUPA PA5400954D, RR-3698-2024):} La Matrícula Regular se aprueba de forma automática en la Intranet DIRCE cuando el horario elegido cumple con los requisitos vigentes; el trámite es gratuito y se atiende en un máximo de dos días hábiles.

    \item \textbf{Regla del Cierre de Matrícula (RR-0570-2022, Art. 78\textdegree{}):} Terminada la primera semana de revisión, la Oficina de Estadística de la Facultad cierra el padrón y envía los archivos digitales a la DIRCE. Después de este cierre ya no se aceptan cambios en la matrícula, salvo con una resolución decanal o rectoral fuera de plazo.

    \item \textbf{Regla del visado obligatorio:} El Reporte y la Boleta de Matrícula solo tienen validez oficial cuando llevan la firma (visado) de la Oficina de Estadística de la Facultad; sin esta firma, no sirven para trámites de becas, seguros o carné.

    \item \textbf{Regla de entrega de listas de clase (RR-0570-2022, Art. 78\textdegree{}):} Las listas actualizadas de alumnos deben entregarse a cada profesor el primer día hábil de la segunda semana de clases.

    \item \textbf{Regla de plazo para reportar a SUNEDU y MINEDU (RR-0570-2022, Art. 79\textdegree{}):} La DIRCE debe enviar los reportes oficiales de la matrícula regular a la SUNEDU y al MINEDU en un máximo de veinte días desde que se genera la matrícula definitiva.

    \item \textbf{Regla de entrega del registro a las Facultades (RR-0570-2022, Art. 80\textdegree{}):} Una vez cerrado el proceso, la DIRCE tiene hasta diez días hábiles para entregar a las Facultades el registro de matrícula, las boletas oficiales y el sistema SIGA actualizado.

    \item \textbf{Regla de autorización del registro digital (RR-0570-2022, Disposición Transitoria Única):} El reglamento autoriza a la DIRCE a simplificar y automatizar los trámites de matrícula, lo que respalda el registro digital que se realiza en esta etapa.
\end{itemize}
